% English solution body for insertion below the existing Task 1 heading.
% Uses standard LaTeX commands only. No additional packages are required.

\subsection*{Conditional independence by d-separation}
An independence statement is implied by the graph if it holds for every
distribution that factorises according to the graph. A path is active given
the observed variables if every non-collider on the path is unobserved and
every collider is observed or has an observed descendant. The graph implies
independence exactly when all paths between the queried nodes are blocked.

\medskip
\noindent
\begin{tabular}{@{}lllp{0.57\linewidth}@{}}
\hline
 & Statement & Implied? & Graphical justification \\
\hline
(a) & $I\perp I$ & No &
A non-degenerate random variable is not independent of itself. The graph
does not require $I$ to be constant. \\[0.5em]
(b) & $B\perp J$ & No &
The directed path $B\to E\to H\to J$ is active. \\[0.5em]
(c) & $F\perp G$ & Yes &
All paths are blocked by unobserved colliders. In particular,
$F\to D\to I\leftarrow E\leftarrow G$ is blocked at $I$. \\[0.5em]
(d) & $C\perp E$ & No &
The common-cause path $C\leftarrow B\to E$ is active. \\[0.5em]
(e) & $A\perp H\mid J$ & No &
The path $A\to C\to D\to I\to J\leftarrow H$ becomes active
when the collider $J$ is observed. \\[0.5em]
(f) & $C\perp E\mid B$ & Yes &
Observing $B$ blocks the common-cause path. Every alternative path
contains an unactivated collider at $I$ or $J$. \\[0.5em]
(g) & $G\perp F\mid I,C$ & No &
Observing $I$ activates the collider on
$G\to E\to I\leftarrow D\leftarrow F$.
The observed node $C$ is outside this path. \\[0.5em]
(h) & $J\perp A\mid C,I$ & No &
The path $A\to C\leftarrow B\to E\to H\to J$ is active:
$C$ is an observed collider, and $I$ is outside the path. \\
\hline
\end{tabular}

\medskip
\noindent\textbf{Final answer: only (c) and (f) are implied by the graph.}
``No'' means that the graph does not guarantee independence. Particular CPDs
can still introduce additional independences.

\subsection*{Detailed graphical reasoning for (h)}
The conditioning set is $Z=\{C,I\}$. Consider the path
\[
 A\longrightarrow C\longleftarrow B\longrightarrow E
 \longrightarrow H\longrightarrow J.
\]
At $C$, both adjacent arrows point into the node, so $C$ is a collider.
Conditioning on $C$ activates this collider. The other internal nodes,
$B$, $E$, and $H$, are non-colliders and none belongs to $Z$.
The observed node $I$ does not lie on this path and therefore does not block it.
Consequently, at least one active path connects $A$ and $J$ given $C,I$.
Thus $J\perp A\mid C,I$ is \textbf{not implied}.

\medskip
\noindent\textit{Note on (a):} The assignment literally states $I\perp I$.
For a discrete variable, independence from itself would require
$P(I=i)=P(I=i)^2$ for every value $i$. Thus every such probability would
have to be either zero or one. The graph imposes no such degeneracy.

\medskip
\noindent\textit{Reference:} UC Berkeley CS188, Section 6.5, D-Separation:
\texttt{https://inst.eecs.berkeley.edu/\string~cs188/textbook/bayes-nets/d-separation.html}.
